% =============================================================================
% Kapitel 2: Theoretische Grundlagen
% =============================================================================

\section{Theoretische Grundlagen}
\label{sec:grundlagen}

\subsection{Verfügbarkeit im Dauerbetrieb}
\label{sec:grundlagen-verfuegbarkeit}

Verfügbarkeit (\textit{Availability}) ist eine zentrale Eigenschaft zuverlässiger
Systeme. Avizienis et al.\ definieren sie als die Bereitschaft eines Systems, seinen
spezifizierten Dienst korrekt zu erbringen \cite{avizienis2004}. Formal ergibt sie sich
aus der mittleren Betriebsdauer zwischen Ausfällen (\textit{Mean Time to Failure}, MTTF)
und der mittleren Wiederherstellungszeit (\textit{Mean Time to Repair}, MTTR):

\begin{equation}
    A = \frac{\text{MTTF}}{\text{MTTF} + \text{MTTR}}
    \label{eq:availability}
\end{equation}

Gleichung~\ref{eq:availability} verdeutlicht, dass hohe Verfügbarkeit nicht allein durch
seltene Ausfälle, sondern ebenso durch kurze Wiederherstellungszeiten erreicht wird. Für
ein System ohne automatische Fehlererkennung ist die MTTR untrennbar an die Zeit
gebunden, bis ein Mensch den Ausfall bemerkt. Während Entwicklungsumgebungen typischerweise
täglich neu gestartet werden und Fehler durch manuelle Beobachtung sichtbar sind,
akkumulieren sich im Dauerbetrieb Effekte, die über kurze Beobachtungsfenster unsichtbar
bleiben. Ein System, das im Entwicklungsbetrieb stabil läuft, kann daher im 24/7-Betrieb
durch eine andere Klasse von Fehlern ausfallen.

\subsection{Fehlerklassifikation}
\label{sec:grundlagen-fehler}

Cristian unterscheidet grundlegende Verhaltensklassen bei Prozessausfällen in verteilten
Systemen \cite{cristian1991}. Für diese Fallstudie sind zwei Klassen relevant. Bei einem
\textit{Fail-Stop}-Ausfall bricht ein fehlerhafter Prozess sofort ab und signalisiert
seinen Ausfall unmissverständlich; der Fehler ist für das Gesamtsystem sichtbar und
adressierbar. Bei einem \textit{Fail-Silent}-Ausfall stellt der Prozess seine Funktion
ohne Meldung ein, das Gesamtsystem läuft scheinbar weiter -- ohne Alarm, ohne Neustart,
ohne Log. Fail-Silent-Verhalten ist in industriellen Systemen besonders gefährlich, weil
der Ausfall nicht durch technische Signale, sondern erst durch ausbleibende
Produktionsergebnisse sichtbar wird.

Davon abzugrenzen sind \textit{schleichende Fehler}, die nicht durch ein singuläres
Ereignis, sondern durch die kontinuierliche Akkumulation eines Ressourcenverbrauchs
entstehen. Charakteristisch ist die erhebliche Zeitspanne zwischen Ursache und sichtbarem
Symptom; im 24/7-Betrieb kann diese Stunden oder ganze Schichten betragen. Birman
beschreibt diesen Zustand als nominal laufendes System, das seinen eigentlichen Zweck
nicht mehr erfüllt \cite{birman2012}.

Ein \textit{Kaskadeneffekt} liegt vor, wenn der Ausfall oder die Verlangsamung einer
Komponente über gemeinsame Schnittstellen den Ausfall weiterer Komponenten verursacht
\cite{ieee_cascade2016}. Solche Effekte sind schwer zu diagnostizieren, da das sichtbare
Symptom nicht die Ursache widerspiegelt: Das System kollabiert an einer anderen Stelle,
als die ursprüngliche Schwachstelle liegt.

\subsection{Prozesssupervision und das Watchdog-Pattern}
\label{sec:grundlagen-supervision}

In langlebigen Multiprocessing-Systemen ist es unzureichend, Prozesse nur beim Start zu
überprüfen. Das \textit{Supervisor-Modell} sieht einen dedizierten Überwachungsprozess
vor, der den Laufzeitstatus aller untergeordneten Prozesse kontinuierlich beobachtet und
bei Ausfall automatisch reagiert -- durch Neustart, Alarm oder geordnetes Herunterfahren.
Dieses Konzept ist im Erlang/OTP-Framework als \textit{Supervision Tree} etabliert, in
dem Supervisoren ihre Kindprozesse überwachen und bei Absturz neu starten
\cite{armstrong2007}. Das \textit{Watchdog-Pattern} ist eine konkrete Ausprägung: Der
Supervisor erwartet in regelmäßigen Abständen ein Lebenszeichen (\textit{Heartbeat}) des
überwachten Prozesses; bleibt es aus, gilt der Prozess als ausgefallen. Für eingebettete
und industrielle Systeme beschreiben Pont und Ong mehrere Varianten dieses Musters mit
unterschiedlichen Reaktionsstrategien \cite{pont_watchdog}.

Fehlt eine solche Logik, entstehen zwei kritische Szenarien. Im ersten hält ein
abgestürzter Prozess, dessen Ausgabe von anderen Prozessen erwartet wird, das System in
einem wartenden Zustand gefangen -- ohne Timeout, ohne Alarm. Im zweiten läuft das System
nach dem Ausfall eines Prozesses scheinbar weiter, Fehler werden erst durch ausbleibende
Ausgaben sichtbar. Beide Szenarien erfordern manuelle Eingriffe, die im 24/7-Betrieb ohne
Bereitschaftsdienst nicht zeitnah erfolgen.

\subsection{Ressourcenerschöpfung und Backpressure}
\label{sec:grundlagen-ressourcen}

Eine \textit{unbounded Queue} wächst ohne Obergrenze; bei anhaltender Produktionsrate und
gedrosseltem Konsum akkumuliert sie Einträge, bis der Speicher erschöpft ist. Eine
\textit{bounded Queue} begrenzt das Wachstum, erzeugt dafür aber \textit{Backpressure}:
Produzenten, die in eine volle Queue schreiben wollen, werden blockiert. Kleppmann
beschreibt Backpressure als notwendigen, aber häufig fehlerhaft implementierten
Mechanismus -- ein blockierendes \texttt{put()} ohne Timeout überträgt den Stau transitiv
auf alle vorgelagerten Komponenten \cite{kleppmann2017}. In einem Multiprocessing-System
mit gemeinsamer Queue kann so ein einzelner langsamer Konsument das gesamte System zum
Einfrieren bringen.

Flüchtige Dateisysteme wie \texttt{tmpfs} und RAM-Disks liegen ausschließlich im
Arbeitsspeicher \cite{linux_tmpfs}. Sie bieten hohe I/O-Geschwindigkeit, sind jedoch auf
eine feste Größe begrenzt: Überschreitet das Schreibvolumen diese Grenze, schlagen
Schreiboperationen mit \texttt{OSError: No space left on device} fehl. Eine RAM-Disk, die
als Puffer zwischen schnellem Schreiben und langsamem Lesen dient, ist strukturell ein
\textit{Single Point of Failure} für alle abhängigen Komponenten.

\subsection{Dev/Prod-Parity und Konfigurationsgefälle}
\label{sec:grundlagen-parity}

Das Prinzip der \textit{Dev/Prod-Parity} (Faktor X des \textit{Twelve-Factor App}-Modells)
fordert, Entwicklungs- und Produktionsumgebung so ähnlich wie möglich zu halten
\cite{wiggins2011}. Abweichungen in Konfigurationsparametern können dazu führen, dass
Fehler in der Entwicklung nicht reproduzierbar sind, die in Produktion deterministisch
auftreten. Ein typisches Muster ist das Konfigurationsgefälle beim Verbindungs-Pooling:
Entwicklungsumgebungen öffnen pro Anfrage eine frische Verbindung, Produktionsumgebungen
nutzen aus Performance-Gründen persistente Verbindungen mit langer Lebensdauer. Die
Fehlerklasse der veralteten Verbindung (\textit{stale connection}) tritt damit im
Entwicklungsbetrieb strukturell nicht auf. Fehlerbehandlungscode kann fehlen, ohne dass
Tests darauf hinweisen; in Produktion führt das erste Auftreten der Fehlerbedingung zum
unkontrollierten Absturz.
